\documentclass[reqno]{amsart}
\usepackage[T1]{fontenc}
\usepackage{amsmath,amssymb,amsthm}
\usepackage{xr-hyper}
\usepackage[hidelinks]{hyperref}
\externaldocument[ps-]{20261125T010500--poincare-series-and-near-orthogonality__workbook}
\externaldocument[dt-]{20261125T002000--definite-theta-lifts__workbook}
\externaldocument[sw-]{20261125T015500--siegel-weil-in-the-smallest-case__workbook}
\externaldocument[bp-]{20261126T020210--bessel-poincare-series-and-their-gram-operators__workbook}
\externaldocument[m1-]{20261126T051900--unfolding-twice-gram-operators-as-sums-over-rational-points__workbook}
\externaldocument[m2-]{20261126T052000--which-terms-survive-characters-minors-and-the-gain-at-the-level__workbook}
\externaldocument[m4-]{20261126T052300--positive-type-and-the-bootstrap-the-second-moment__workbook}

% Chapter number; \ref{chap} gives it.
\newcounter{chap}
% Chapter numbers, generated from 20261125T012000--grh-workbook__guide.org
% by scripts/make-chapters.py.  Do not edit: change the order of the
% index and rerun the script (./preview-site.sh and CI do this).
%
% Usage in a sheet, after \newcounter{chap}:
%   % Chapter numbers, generated from 20261125T012000--grh-workbook__guide.org
% by scripts/make-chapters.py.  Do not edit: change the order of the
% index and rerun the script (./preview-site.sh and CI do this).
%
% Usage in a sheet, after \newcounter{chap}:
%   % Chapter numbers, generated from 20261125T012000--grh-workbook__guide.org
% by scripts/make-chapters.py.  Do not edit: change the order of the
% index and rerun the script (./preview-site.sh and CI do this).
%
% Usage in a sheet, after \newcounter{chap}:
%   \input{chapters}\chapnumber{<sheet ID>}
\makeatletter
\providecommand\chapters@set[2]{\@namedef{chapters@#1}{#2}}
\chapters@set{20261125T001000}{1}% zeros-to-ramp-sums
\chapters@set{20261125T001100}{2}% inert-primes-and-estimate-b
\chapters@set{20261125T005000}{3}% galois-fourier-and-averaging
\chapters@set{20261125T005500}{4}% columns-from-a-sphere
\chapters@set{20261125T000000}{5}% slope-method
\chapters@set{20261126T030100}{6}% a-twist-for-each-prime
\chapters@set{20261125T003000}{7}% the-chain-and-the-square-root
\chapters@set{20261125T010000}{8}% rank-lemma-and-zero-estimate
\chapters@set{20261125T010500}{9}% poincare-series-and-near-orthogonality
\chapters@set{20261124T223000}{10}% elliptic-curves-for-the-proof
\chapters@set{20261125T009500}{11}% fock-model-and-theta-intertwiners
\chapters@set{20261124T223100}{12}% complex-multiplication-and-inert-reduction
\chapters@set{20261125T008000}{13}% one-cm-determinant
\chapters@set{20261124T214500}{14}% supersingular-theta-divisibility
\chapters@set{20261124T214600}{15}% the-isotropic-plane
\chapters@set{20261125T013700}{16}% pairings-on-a-formal-group-of-height-two
\chapters@set{20261125T013500}{17}% hecke-nondegeneracy-on-the-tree
\chapters@set{20261126T030200}{18}% the-bound-is-attained
\chapters@set{20261125T009000}{19}% kronecker-limit-and-chowla-selberg
\chapters@set{20261125T002000}{20}% definite-theta-lifts
\chapters@set{20261125T015500}{21}% siegel-weil-in-the-smallest-case
\chapters@set{20261125T002100}{22}% definite-theta-lifts-and-norms
\chapters@set{20261125T016500}{23}% hecke-operators-and-theta-lifts
\chapters@set{20261125T016700}{24}% divisibility-along-the-chain
\chapters@set{20261125T004000}{25}% denominators-and-sha
\chapters@set{20261125T004100}{26}% tate-twists-and-sha
\chapters@set{20261126T020430}{27}% the-minus-part-by-gauss-sums
\chapters@set{20261125T014100}{28}% congruence-modules
\chapters@set{20261125T014500}{29}% eisenstein-congruences-and-selmer-classes
\chapters@set{20261125T014700}{30}% pseudocharacters-and-the-length-bound
\chapters@set{20261125T011000}{31}% galois-representations-black-box
\chapters@set{20261126T030300}{32}% the-data
\chapters@set{20261126T020210}{33}% bessel-poincare-series-and-their-gram-operators
\chapters@set{20261126T030400}{34}% the-seesaw-identity
\chapters@set{20261126T030500}{35}% divisibility-of-the-composite-lift
\chapters@set{20261126T030600}{36}% theorem-a-short-forms
\chapters@set{20261126T030700}{37}% the-determinant-of-a-chain
\chapters@set{20261126T030800}{38}% the-cost-of-one-lift
\chapters@set{20261126T030900}{39}% denominators-of-projections-to-lifts
\chapters@set{20261126T031000}{40}% twists-with-few-congruences
\chapters@set{20261126T031100}{41}% theorem-c-and-the-main-inequality
\providecommand\chapnumber[1]{%
  \@ifundefined{chapters@#1}%
    {\PackageError{chapters}{Sheet #1 is not listed in the index}%
      {List it in 20261125T012000--grh-workbook__guide.org\MessageBreak
       and run scripts/make-chapters.py.}}%
    {\edef\thechap{\@nameuse{chapters@#1}}}}
\makeatother
\chapnumber{<sheet ID>}
\makeatletter
\providecommand\chapters@set[2]{\@namedef{chapters@#1}{#2}}
\chapters@set{20261125T001000}{1}% zeros-to-ramp-sums
\chapters@set{20261125T001100}{2}% inert-primes-and-estimate-b
\chapters@set{20261125T005000}{3}% galois-fourier-and-averaging
\chapters@set{20261125T005500}{4}% columns-from-a-sphere
\chapters@set{20261125T000000}{5}% slope-method
\chapters@set{20261126T030100}{6}% a-twist-for-each-prime
\chapters@set{20261125T003000}{7}% the-chain-and-the-square-root
\chapters@set{20261125T010000}{8}% rank-lemma-and-zero-estimate
\chapters@set{20261125T010500}{9}% poincare-series-and-near-orthogonality
\chapters@set{20261124T223000}{10}% elliptic-curves-for-the-proof
\chapters@set{20261125T009500}{11}% fock-model-and-theta-intertwiners
\chapters@set{20261124T223100}{12}% complex-multiplication-and-inert-reduction
\chapters@set{20261125T008000}{13}% one-cm-determinant
\chapters@set{20261124T214500}{14}% supersingular-theta-divisibility
\chapters@set{20261124T214600}{15}% the-isotropic-plane
\chapters@set{20261125T013700}{16}% pairings-on-a-formal-group-of-height-two
\chapters@set{20261125T013500}{17}% hecke-nondegeneracy-on-the-tree
\chapters@set{20261126T030200}{18}% the-bound-is-attained
\chapters@set{20261125T009000}{19}% kronecker-limit-and-chowla-selberg
\chapters@set{20261125T002000}{20}% definite-theta-lifts
\chapters@set{20261125T015500}{21}% siegel-weil-in-the-smallest-case
\chapters@set{20261125T002100}{22}% definite-theta-lifts-and-norms
\chapters@set{20261125T016500}{23}% hecke-operators-and-theta-lifts
\chapters@set{20261125T016700}{24}% divisibility-along-the-chain
\chapters@set{20261125T004000}{25}% denominators-and-sha
\chapters@set{20261125T004100}{26}% tate-twists-and-sha
\chapters@set{20261126T020430}{27}% the-minus-part-by-gauss-sums
\chapters@set{20261125T014100}{28}% congruence-modules
\chapters@set{20261125T014500}{29}% eisenstein-congruences-and-selmer-classes
\chapters@set{20261125T014700}{30}% pseudocharacters-and-the-length-bound
\chapters@set{20261125T011000}{31}% galois-representations-black-box
\chapters@set{20261126T030300}{32}% the-data
\chapters@set{20261126T020210}{33}% bessel-poincare-series-and-their-gram-operators
\chapters@set{20261126T030400}{34}% the-seesaw-identity
\chapters@set{20261126T030500}{35}% divisibility-of-the-composite-lift
\chapters@set{20261126T030600}{36}% theorem-a-short-forms
\chapters@set{20261126T030700}{37}% the-determinant-of-a-chain
\chapters@set{20261126T030800}{38}% the-cost-of-one-lift
\chapters@set{20261126T030900}{39}% denominators-of-projections-to-lifts
\chapters@set{20261126T031000}{40}% twists-with-few-congruences
\chapters@set{20261126T031100}{41}% theorem-c-and-the-main-inequality
\providecommand\chapnumber[1]{%
  \@ifundefined{chapters@#1}%
    {\PackageError{chapters}{Sheet #1 is not listed in the index}%
      {List it in 20261125T012000--grh-workbook__guide.org\MessageBreak
       and run scripts/make-chapters.py.}}%
    {\edef\thechap{\@nameuse{chapters@#1}}}}
\makeatother
\chapnumber{<sheet ID>}
\makeatletter
\providecommand\chapters@set[2]{\@namedef{chapters@#1}{#2}}
\chapters@set{20261125T001000}{1}% zeros-to-ramp-sums
\chapters@set{20261125T001100}{2}% inert-primes-and-estimate-b
\chapters@set{20261125T005000}{3}% galois-fourier-and-averaging
\chapters@set{20261125T005500}{4}% columns-from-a-sphere
\chapters@set{20261125T000000}{5}% slope-method
\chapters@set{20261126T030100}{6}% a-twist-for-each-prime
\chapters@set{20261125T003000}{7}% the-chain-and-the-square-root
\chapters@set{20261125T010000}{8}% rank-lemma-and-zero-estimate
\chapters@set{20261125T010500}{9}% poincare-series-and-near-orthogonality
\chapters@set{20261124T223000}{10}% elliptic-curves-for-the-proof
\chapters@set{20261125T009500}{11}% fock-model-and-theta-intertwiners
\chapters@set{20261124T223100}{12}% complex-multiplication-and-inert-reduction
\chapters@set{20261125T008000}{13}% one-cm-determinant
\chapters@set{20261124T214500}{14}% supersingular-theta-divisibility
\chapters@set{20261124T214600}{15}% the-isotropic-plane
\chapters@set{20261125T013700}{16}% pairings-on-a-formal-group-of-height-two
\chapters@set{20261125T013500}{17}% hecke-nondegeneracy-on-the-tree
\chapters@set{20261126T030200}{18}% the-bound-is-attained
\chapters@set{20261125T009000}{19}% kronecker-limit-and-chowla-selberg
\chapters@set{20261125T002000}{20}% definite-theta-lifts
\chapters@set{20261125T015500}{21}% siegel-weil-in-the-smallest-case
\chapters@set{20261125T002100}{22}% definite-theta-lifts-and-norms
\chapters@set{20261125T016500}{23}% hecke-operators-and-theta-lifts
\chapters@set{20261125T016700}{24}% divisibility-along-the-chain
\chapters@set{20261125T004000}{25}% denominators-and-sha
\chapters@set{20261125T004100}{26}% tate-twists-and-sha
\chapters@set{20261126T020430}{27}% the-minus-part-by-gauss-sums
\chapters@set{20261125T014100}{28}% congruence-modules
\chapters@set{20261125T014500}{29}% eisenstein-congruences-and-selmer-classes
\chapters@set{20261125T014700}{30}% pseudocharacters-and-the-length-bound
\chapters@set{20261125T011000}{31}% galois-representations-black-box
\chapters@set{20261126T030300}{32}% the-data
\chapters@set{20261126T020210}{33}% bessel-poincare-series-and-their-gram-operators
\chapters@set{20261126T030400}{34}% the-seesaw-identity
\chapters@set{20261126T030500}{35}% divisibility-of-the-composite-lift
\chapters@set{20261126T030600}{36}% theorem-a-short-forms
\chapters@set{20261126T030700}{37}% the-determinant-of-a-chain
\chapters@set{20261126T030800}{38}% the-cost-of-one-lift
\chapters@set{20261126T030900}{39}% denominators-of-projections-to-lifts
\chapters@set{20261126T031000}{40}% twists-with-few-congruences
\chapters@set{20261126T031100}{41}% theorem-c-and-the-main-inequality
\providecommand\chapnumber[1]{%
  \@ifundefined{chapters@#1}%
    {\PackageError{chapters}{Sheet #1 is not listed in the index}%
      {List it in 20261125T012000--grh-workbook__guide.org\MessageBreak
       and run scripts/make-chapters.py.}}%
    {\edef\thechap{\@nameuse{chapters@#1}}}}
\makeatother

\ifcsname chapters@20261126T052200\endcsname
  \chapnumber{20261126T052200}
\else
  \renewcommand{\thechap}{?}
\fi
\renewcommand{\thesection}{\thechap.\arabic{section}}

\theoremstyle{definition}
\newtheorem{problem}{Problem}[section]
\newtheorem{definition}{Definition}[section]
\theoremstyle{plain}
\newtheorem{theorem}{Theorem}[section]
\theoremstyle{remark}
\newtheorem*{leanremark}{In the Lean formalization}

\newcommand{\Z}{\mathbb{Z}}
\newcommand{\Q}{\mathbb{Q}}
\newcommand{\R}{\mathbb{R}}
\newcommand{\C}{\mathbb{C}}
\newcommand{\A}{\mathbb{A}}
\newcommand{\ii}{\mathrm{i}}
\newcommand{\Uc}{\mathcal{U}}
\newcommand{\Dc}{\mathcal{D}}
\newcommand{\cMc}{\mathcal{M}}
\newcommand{\Zc}{\mathcal{Z}}
\newcommand{\Oh}{\hat{\mathcal O}_E}
\newcommand{\hd}{\operatorname{hd}}
\newcommand{\lE}{\operatorname{le}}
\newcommand{\disc}{\operatorname{disc}}
\newcommand{\vol}{\operatorname{vol}}
\newcommand{\Dm}{D_{\mathrm{main}}}
\newcommand{\Di}{D_{\mathrm{id}}}
\newcommand{\tl}[1]{\texttt{#1}}
\newcommand{\glsrc}[3]{\href{https://github.com/lsernau/grh-lean/blob/72fead1/grh/src/#1\#L#2}{#3}}

\setlength{\emergencystretch}{3em}

\title{The open cell: Siegel's formula and Poisson summation}
\author{}
\date{November 26, 2026}

\begin{document}
\maketitle
\refstepcounter{chap}\label{chap}

\paragraph{Prerequisites.}
From Chapter~\ref{bp-chap}: the hermitian space \(Y_0=\bigoplus_aEv_a\oplus X_0\oplus Ez\oplus\bigoplus_aEv'_a\) with \(h(z,z)=-N\), the unipotent group \(\Uc\) (Section~\ref{bp-sec:subgroup}), and the covolume \(\nu\) (Definition~\ref{bp-def:scaled}).  From Chapter~\ref{m1-chap}: the degrees \(d_\tau\) of Theorem~\ref{m1-thm:bessel-system}, \(\deg=\prod_\tau d_\tau\), and the sizes \(\Di\), \(\Dm=D_X^2/(\deg\cdot\mathrm{mass}_r)\) of Definition~\ref{m1-def:two-sizes}.  From Chapter~\ref{m2-chap}: the frame and the leading minors \(\langle\gamma\rangle_j\) (Definition~\ref{m2-def:frame}), \(\disc\) (Definition~\ref{m2-def:disc}), links, open reductions, fixers and link second moments (Definitions~\ref{m2-def:link}, \ref{m2-def:open}, \ref{m2-def:fixer}, \ref{m2-def:moments}).  From Chapter~\ref{sw-chap}: the densities \(\alpha_\infty\), \(\alpha_p\) of Section~\ref{sw-sec:densities} and Braun's formula~\eqref{sw-eq:sw-herm}.  Problem~\ref{dt-pr:class-one}; Problem~\ref{ps-pr:formal-degree}.  Fourier analysis on finite abelian groups; Hensel's lemma; Schur orthogonality for square-integrable representations \cite[Ch.~IX]{Kn}.  Hermitian forms are linear in the first variable, and \(\lvert\cdot\rvert^2\) is the product with the complex conjugate.  In this sheet \(j\), \(l\) are integers or residues, \(y\) is a vector, \(\omega\) is an isometry, \(\mathrm e(\cdot):=e^{2\pi\ii\,\cdot}\) on \(\R\), and \(K'\) is a level (not the field \(K\)).

\medskip
The second moment of the Gram operators of Chapter~\ref{bp-chap} is bounded in Theorem~\ref{m4-thm:second-moment} by splitting the reductions \(s\in\bar G\) into the open ones and the rest.  This sheet treats the open ones.  Over the open reductions that fix chosen vectors of the span of the frame at \(\#\Dc\) chosen places, the link second moment of the terms with \(\disc\ne0\) is at most \((1+b^{-4})\,b^{-6k\#\Dc}\,\Dm\,\nu^2\) (Theorem~\ref{thm:open}).  Where does \(\Dm\) come from?  A term over an open reduction is a head, a last entry and a residual isometry.  The residual isometries with a given last entry are counted by vectors of a given norm, and Siegel's formula counts these exactly.  The heads are summed by Poisson summation, one entry at a time: the zero frequencies give an integral, and the other frequencies are small.  The bound has no small instance: in the proof \(\varrho\ge10^5\).  We meet the steps separately: one isometry over \(\Q\), six squares, and the groups \(\Z/5\) and \(\Z/3\).

\section{Head, last entry and residual isometry}\label{sec:last}

Let \(V=\bigoplus_{a\le\varrho}Ev_a\oplus X\oplus Ez\oplus\bigoplus_{a\le\varrho}Ev'_a\) be a hermitian space of the shape of Section~\ref{bp-sec:subgroup}, with \(X\) of rank \(k\) and totally positive definite, and \(h(z,z)=-N\).  A \emph{residual isometry} is an isometry \(\omega\) of \(V\) with \(\omega v_a=v_a\) and \(\omega v'_a=v'_a\) for all \(a\).  It preserves the orthogonal complement \(X\oplus Ez\) of the \(v_a\), \(v'_a\).

\begin{problem}\label{pr:last-entry}
Let \(\omega\) be a residual isometry, and write \(\omega z=y+c_\omega z\) with \(y\in X\), \(c_\omega\in E\).  Then \(\lvert h(z,\omega z)\rvert^2=N^2\lvert c_\omega\rvert^2\) and
\[
 N\,h(y,y)=\lvert h(z,\omega z)\rvert^2-N^2 .
\]
Hence \(\lvert\tau h(z,\omega z)\rvert\ge N\) at every infinite place \(\tau\), with strict inequality at every \(\tau\) if \(y\ne0\) and equality at every \(\tau\) if \(y=0\).
\end{problem}

Now let \(V=Y_0\), with the frame \((v_1,\dots,v_\varrho,z)\) and the leading minors \(\langle\gamma\rangle_j\), \(0\le j\le\varrho+1\), of Definition~\ref{m2-def:frame}, so that \(\langle\gamma\rangle_0=1\) and \(\langle\gamma\rangle_{\varrho+1}\) is the determinant of the matrix of frame products.

\begin{definition}\label{def:head}
Let \(\gamma\in U(Y_0)(F)\) have \(\langle\gamma\rangle_j\ne0\) for \(1\le j\le\varrho\).  Its \emph{head} is \(\hd(\gamma):=\bigl(\langle\gamma\rangle_{a+1}/\langle\gamma\rangle_a\bigr)_{0\le a<\varrho}\in(E^\times)^\varrho\).  If also \(\langle\gamma\rangle_{\varrho+1}\ne0\), its \emph{last entry} is \(\lE(\gamma):=\langle\gamma\rangle_{\varrho+1}/\langle\gamma\rangle_\varrho\).  A \emph{residual isometry of the last floor} is a residual isometry of \(Y_0\).
\end{definition}

By Definition~\ref{m2-def:disc}, \(\disc(\gamma)=\lvert\langle\gamma\rangle_\varrho\rvert^2\bigl(\lvert\lE(\gamma)\rvert^2-N^2\bigr)\) whenever \(\lE(\gamma)\) is defined.  In the proof, if \(s\) is open, then every term \(\gamma\) over \((c,s,c')\) has \(\langle\gamma\rangle_j\ne0\) for \(j\le\varrho\).  Moreover \(\gamma=n^{-1}\mathfrak w\,\omega\,n'\) with \(n,n'\in\Uc(F)\), \(\omega\) a residual isometry of \(Y_0\), and \(\mathfrak w\) a Weyl element that exchanges each line \(Ev_a\) with \(Ev'_a\), fixes \(X_0\oplus Ez\) and depends only on \(\hd(\gamma)\).  The \((z,z)\) entry \(h(\omega z,z)\) of the matrix of frame products of \(\omega\) is \(\lE(\gamma)\), so \(\lvert\lE(\gamma)\rvert=\lvert h(z,\omega z)\rvert\).  By Problem~\ref{pr:last-entry} the last entry is therefore always defined there, and \(\disc(\gamma)\ne0\) exactly when the \(X_0\)-part \(y\) of \(\omega z\) is nonzero.

Consider a block of the link second moment over an open reduction \(s\) whose sum is nonzero.  In the proof one residual isometry \(\omega\) names the link, and the heads of the contributing terms lie in a finite set \(\mathcal E^\varrho\), where \(\mathcal E\subset E^\times\) consists of quotients of integers of bounded size.  The block is the sum, over \(\mathfrak h\in\mathcal E^\varrho\), of the sums of the amplitudes of the \(n^{-1}\mathfrak w\,\omega\,n'\) over \((c,s,c')\), for \(n,n'\in\Uc(F)\), where \(\mathfrak w\) is the Weyl element of \(\mathfrak h\).  Different nonzero blocks give different \((c,c',s,\omega)\).  So the heads are summed inside the square, and the residual isometries outside it.  By Problem~\ref{pr:last-entry}, the sum over the \(\omega\) with a given last entry is controlled by a sum over vectors \(y\in X_0\) with \(N\,h(y,y)=\lvert\lE\rvert^2-N^2\), one for each class.

\section{Siegel's formula for six squares}\label{sec:siegel}

The smallest rank to which Braun's formula~\eqref{sw-eq:sw-herm} applies is \(3\), and for \(\Z[\ii]^3\) the class set is \(\{I_3\}\) (Problem~\ref{dt-pr:class-one}).  So the mean over the class set is a count: \(r_{I_3}(j):=\#\{y\in\Z[\ii]^3:y^*y=j\}\) equals \(\alpha_\infty(j)\prod_p\alpha_p(j)\), with the densities of Section~\ref{sw-sec:densities} for the form \(Q(y)=y^*y\), a sum of six squares on \(\Z^6\).  Let \(\chi_{-4}\) be the nontrivial character modulo \(4\).

\begin{problem}\label{pr:six-squares}
By machine, \(r_{I_3}(1)=12\), \(r_{I_3}(2)=60\) and \(r_{I_3}(5)=312\).
\end{problem}

\begin{problem}\label{pr:local-counts}
For odd \(p\), the number of \(x\in(\Z/p)^6\) with \(\sum_ix_i^2=1\) is \(p^5-\chi_{-4}(p)p^2\), and the same holds with \(1\) replaced by any \(j\) prime to \(p\): by machine \(252\), \(3100\), \(16856\) for \(p=3,5,7\).  By machine, there are \(61236=3^5\cdot252\) solutions modulo \(9\), and \(3\cdot2^{5e}/4\) solutions modulo \(2^e\) for \(e=2,3\) (\(768\), \(24576\)).  For \(e\ge3\) the count modulo \(2^{e+1}\) is \(2^5\) times the count modulo \(2^e\).  Hence \(\alpha_p(1)=1-\chi_{-4}(p)p^{-3}\) for odd \(p\), and \(\alpha_2(1)=3/4\).
\end{problem}

For \(p\equiv3\pmod4\), \(y\mapsto y^*y\) modulo \(p\) is the hermitian form \(\sum_iy_i^{p+1}\) on \(\mathbb F_{p^2}^3\).

\begin{problem}\label{pr:isotropic}
Let \(q\) be a prime power.  The form \(\sum_{i\le n}x_i^{q+1}\) on \(\mathbb F_{q^2}^n\) has \((q^n-(-1)^n)(q^{n-1}-(-1)^{n-1})\) nonzero isotropic vectors: \(9\), \(27\), \(135\) for \(q=2\) and \(n=2,3,4\).  Divided by \(q^{2n-1}\) this is \((1-(-1)^nq^{-n})(1+(-1)^nq^{1-n})\).
\end{problem}

\begin{problem}\label{pr:real-density}
For real \(j>0\), \(\vol\{y\in\C^3:\lvert y\rvert^2\le j\}=\pi^3j^3/6\), so \(\alpha_\infty(j)=\pi^3j^2/2\).
\end{problem}

\begin{problem}\label{pr:densities}
With \(L(3,\chi_{-4})=\pi^3/32\): \(\frac{\pi^3}2\cdot\frac34\cdot\prod_{p\text{ odd}}\bigl(1-\chi_{-4}(p)p^{-3}\bigr)=12=r_{I_3}(1)\).  For \(j=5\) the factor at \(5\) is \(\alpha_5(5)=3224/3125\), the factor at infinity is \(\alpha_\infty(5)=25\pi^3/2\), the other factors are those of \(j=1\), and the product is \(312=r_{I_3}(5)\).
\end{problem}

The general statement is for a CM field with Weil data: \(E\) contains a primitive \(5040\)th root of unity, and there are a \(j_E\in E\) with \(\bar j_E=-j_E\), \(\operatorname{Im}\tau(j_E)>0\) at every \(\tau\) and \((1-j_E^{-1})/2\) an algebraic integer, a finite idele that makes \(\Oh\) self-dual for the pairing defined by \(j_E\) and the additive character \(\psi_f\) of the finite adeles, and a splitting character that is \(-1\) at the chosen uniformizer of every finite place fixed by complex conjugation \cite{GL}.  Here \(\vol(\Oh)=2^f\lvert d_E\rvert^{-1/2}\), with \(d_E\) the discriminant of \(E\).

\begin{theorem}[Siegel's formula for one vector]\label{thm:siegel}
Let \(E\) be a CM field with Weil data, \(k\ge3\), \(K'\) an open subgroup of the stabilizer of \(\Oh^k\) in the finite adelic unitary group of \(y^*y\), \(\phi\ge0\) a locally constant, compactly supported, \(K'\)-invariant function on \(\A_{E,f}^k\), and \(\mathfrak n\in E\) with \(\bar{\mathfrak n}=\mathfrak n\) and \(\tau\mathfrak n>0\) at every infinite place.  Then each sum \(\sum_{y\in E^k,\,y^*y=\mathfrak n}\phi(g^{-1}y)\) converges, and
\begin{multline*}
 \sum_{[g]}\frac{1}{\#\mathrm{Aut}(g)}\sum_{\substack{y\in E^k\\ y^*y=\mathfrak n}}\phi(g^{-1}y)\\
 =\Bigl(\sum_{[g]}\frac{1}{\#\mathrm{Aut}(g)}\Bigr)\lvert d_E\rvert^{1/4}\prod_\tau\frac{\pi^k(\tau\mathfrak n)^{k-1}}{(k-1)!}
 \int\!\!\int\phi(u)\,\psi_f\bigl(\tfrac12\zeta(u^*u-\mathfrak n)\bigr)\,du\,d\zeta ,
\end{multline*}
where \([g]\) runs over the classes of level \(K'\), \(u\) over \(\A_{E,f}^k\) with \(\Oh\) of volume \(\vol(\Oh)\) in each coordinate, and \(\zeta\) over the self-adjoint finite adeles with the codifferent of volume \(1\).
\end{theorem}

\noindent\emph{Status}: classical; due to Siegel \cite{Si} for quadratic forms and Braun \cite{Br} for hermitian forms, and a case of Weil's theorem \cite[Thm.~5]{We65}.  The last integral is the product of the local densities, left as one adelic integral; its normalization is not compared here with that of Chapter~\ref{sw-chap}.  If there is one class, or every class has trivial automorphism group, the formula is a statement about counts; Problem~\ref{pr:six-squares} is the case of one class, over a field without Weil data.

In the proof, Theorem~\ref{thm:siegel} is applied, for each last entry, to the vectors \(y\in X_0\) with \(N\,h(y,y)=\lvert\lE\rvert^2-N^2\), that is at \(\mathfrak n=(\lE\,\overline{\lE}-N^2)/N\).  It is used as an upper bound, with an extra factor \(1+b^{-6}\) from the local densities.

\section{Poisson summation, one entry at a time}\label{sec:poisson}

In the proof, each double-coset sum of Section~\ref{sec:last} is a Kloosterman sum at finitely many finite places times an orbital integral at the infinite places.  Squaring the sum over heads gives a sum over pairs of heads, and the pairs of different heads are kept.  They are summed by Poisson summation in one entry at a time.  For a function \(\mathfrak f\) on \(\Z/p\) let \(\hat{\mathfrak f}(l):=p^{-1}\sum_j\mathfrak f(j)\,\mathrm e(-lj/p)\), so that \(\mathfrak f(j)=\sum_l\hat{\mathfrak f}(l)\,\mathrm e(lj/p)\).

\begin{problem}\label{pr:fourier-five}
On \(\Z/5\), \(\mathfrak f(j)=2+\cos(2\pi j/5)\) has \(\hat{\mathfrak f}=(2,\tfrac12,0,0,\tfrac12)\).  So \(\mathfrak f(0)=3\), while the zero frequency alone gives \(2\).
\end{problem}

\begin{problem}\label{pr:poisson-finite}
Let \(\mathfrak f\) be a function on a finite abelian group \(\mathcal G\), and \(\hat{\mathfrak f}(\psi'):=\#\mathcal G^{-1}\sum_{j\in\mathcal G}\mathfrak f(j)\overline{\psi'(j)}\) for the characters \(\psi'\) of \(\mathcal G\).  Then \(\lvert\mathfrak f(j)\rvert\le\sum_{\psi'}\lvert\hat{\mathfrak f}(\psi')\rvert\) for every \(j\), and \(\hat{\mathfrak f}(\psi_0)\) is the mean of \(\mathfrak f\), for the trivial character \(\psi_0\).  If \(\mathfrak f\) is real and \(\sum_{\psi'\ne\psi_0}\lvert\hat{\mathfrak f}(\psi')\rvert\le c\), then \(\mathfrak f(j)\le\hat{\mathfrak f}(\psi_0)+c\).
\end{problem}

For \(j\in\Z/p\) let \(\mathrm{Kl}(j):=\sum_{l\in(\Z/p)^\times}\mathrm e\bigl((l+jl^{-1})/p\bigr)\), a Kloosterman sum.

\begin{problem}\label{pr:kloosterman}
For \(p=5\) the values \(\mathrm{Kl}(0),\dots,\mathrm{Kl}(4)\) are
\[
 -1,\quad \tfrac{3-\sqrt5}2,\quad -1-\sqrt5,\quad \sqrt5-1,\quad \tfrac{3+\sqrt5}2,
\]
and \(\sum_j\lvert\mathrm{Kl}(j)\rvert^2=20\).
\end{problem}

\begin{problem}\label{pr:kloosterman-sum}
\(\sum_{j\in\Z/p}\lvert\mathrm{Kl}(j)\rvert^2=p(p-1)\).  In the expansion of the square, the terms with \(l\ne l'\) cancel after the sum over \(j\): the matrix \(\bigl(\sum_j\mathrm e(j(l^{-1}-l'^{-1})/p)\bigr)_{l,l'}\) is \(p\) times the identity.
\end{problem}

So once a free variable is summed, different heads are orthogonal.  At a fixed nonzero frequency they are not orthogonal, but the frequency gains a factor in each free direction.

\begin{problem}\label{pr:three}
Let \(\mathfrak A_l:=\bigl(\mathrm e(ljj'/3)\bigr)_{j,j'\in\Z/3}\).  Then \(\mathfrak A_1\mathfrak A_1^*=3\cdot1\), and \(\mathfrak A_0\mathfrak A_0^*\) is the matrix with all entries \(3\).  Their squared operator norms are \(3\) and \(9\), and those of the \(k\)-th tensor powers are \(3^k\) and \(9^k\).
\end{problem}

In the proof, the rank \(k\) of \(X_0\) pays for the nonzero frequencies in the same way.  At each infinite place the bound at a nonzero frequency decays at least like the \((k-(\varrho+1)^7)\)-th power of its inverse size.  We do not derive these local bounds here.

\begin{problem}\label{pr:schur}
The entrywise product of two positive semidefinite matrices of the same size is positive semidefinite.
\end{problem}

In the proof, the form on pairs of heads with a common last entry is the entrywise product of positive semidefinite kernels, one at each infinite place and one at the finite places.  So it is positive semidefinite, by Problem~\ref{pr:schur}, and the densities obtained from it by integrating out entries are nonnegative.  Nonnegativity is what Theorem~\ref{thm:poisson} below needs.

\begin{problem}\label{pr:formal-degree-sum}
Let \(\pi\) be an irreducible unitary representation of a unimodular group \(\mathcal G'\) with square-integrable matrix coefficients and formal degree \(d_\pi\), so that \(\int_{\mathcal G'}\langle\pi(g)u_1,u_3\rangle\overline{\langle\pi(g)u_2,u_4\rangle}\,dg=d_\pi^{-1}\langle u_1,u_2\rangle\overline{\langle u_3,u_4\rangle}\).  For every vector \(u\) with \(\lVert u\rVert^2=D'\),
\[
 \int_{\mathcal G'}\lvert\langle\pi(g)u,u\rangle\rvert^2\,dg=\frac{D'^2}{d_\pi} .
\]
For the representation of Problem~\ref{ps-pr:formal-degree} and \(D'=1\), this is the reciprocal of the formal degree found there.
\end{problem}

Identity (c) of Theorem~\ref{m1-thm:bessel-system} at each infinite place \(\tau\) is the relation of Problem~\ref{pr:formal-degree-sum} for the single vector \(W_\tau\), whose squared norm is \(\#I_\tau\), with \(d_\pi=d_\tau\).  In the proof, at the zero frequency the integral of a form over a common entry is exactly the form of one length less, and the form of length zero has modulus at most \((1+b^{-10})/d_\tau\) at the infinite place \(\tau\).  This is the factor \(\prod_\tau1/d_\tau=1/\deg\) of \(c_{\mathrm{bot}}\) below.  Poisson summation is used in the following one-sided form.

\begin{theorem}[one-sided Poisson summation]\label{thm:poisson}
Let \(E\) be a totally complex number field with chosen complex embeddings, \(\psi\) the standard character of \(\A_E/E\), and \(du\) the Haar measure on \(\A_E\) for which \(E\) has covolume \(1\).  Let \(\mathfrak f\colon\A_E\to\R\) be nonnegative and integrable, and lower semicontinuous at each point of a finite set \(\Zc\subset E\), and put \(\hat{\mathfrak f}(l):=\int\mathfrak f(u)\psi(-lu)\,du\).  If \(\sum_{l\in\Zc'}\lvert\hat{\mathfrak f}(l)\rvert\le c_{\mathfrak f}\) for every finite \(\Zc'\subset E\), then \(\sum_{u\in\Zc}\mathfrak f(u)\le c_{\mathfrak f}\).
\end{theorem}

\noindent\emph{Status}: an inequality form of Poisson summation on \(\A_E/E\) \cite{Ta}, for functions that need not be continuous.  We use it as the proof proves it, where the test function is smoothed by a bump on the compact quotient.

\begin{theorem}[the sieve]\label{thm:sieve}
In the proof, fix \(\Dc\) and the vectors of Theorem~\ref{thm:open}.  There are real numbers \(\Sigma_0,\dots,\Sigma_{\varrho+1}\), the sums of the densities of the pairs of heads of length \(i\) (tuples of \(i\) entries) with entries in \(\mathcal E\) and a common last entry (for \(i=0\), the density of the empty pair), and \(\lambda_0,\dots,\lambda_\varrho\ge0\), the sums of the weights of the heads of length \(i\) with entries in \(\mathcal E\).  They satisfy, for \(0\le i\le\varrho\),
\[
 \Sigma_0\le c_{\mathrm{bot}},\qquad \Sigma_{i+1}\le c_i\bigl(\Sigma_i+2c_{\mathrm{fr}}\lambda_i^2\bigr),\qquad \lambda_i^2\le c_{\mathrm{wt}}\prod_{i'<i}c_{i'} ,
\]
and \(\Sigma_{\varrho+1}\) bounds the sum, over the last entries, of the top form over all pairs of heads.  Here \(c_i>0\) is the constant of Poisson summation in the entry \(i\): \(\vol(\Oh)\) divided by the Jacobian, at the infinite places, of the rescaling that the proof applies to that entry.  Further \(c_{\mathrm{bot}}=\prod_\tau\bigl((1+b^{-10})/d_\tau\bigr)\cdot(1+b^{-6})\,b^{-6k\#\Dc}\), \(c_{\mathrm{fr}}=b^{2002(\varrho+1)^9f}c_{\mathrm{dec}}\) with an explicit constant \(c_{\mathrm{dec}}>0\), and \(c_{\mathrm{wt}}=b^{400(\varrho+1)^9f}\).
\end{theorem}

\noindent\emph{Status}: read from the proof.  Each step of the recurrence is Theorem~\ref{thm:poisson} in the entry \(i\).  The zero frequency gives \(c_i\) times the sum of the densities over all pairs of heads of length \(i\); the pairs with a common last entry give \(c_i\Sigma_i\), and those whose last entries differ give the error \(c_ic_{\mathrm{fr}}\lambda_i^2\).  The nonzero frequencies give the other error \(c_ic_{\mathrm{fr}}\lambda_i^2\).  The local bounds behind \(c_{\mathrm{fr}}\) are not derived here.

\begin{problem}\label{pr:induction}
Under Theorem~\ref{thm:sieve}, \(\Sigma_i\le\bigl(\prod_{i'<i}c_{i'}\bigr)\bigl(c_{\mathrm{bot}}+2i\,c_{\mathrm{wt}}c_{\mathrm{fr}}\bigr)\) for \(0\le i\le\varrho+1\).
\end{problem}

In the proof, Theorem~\ref{thm:siegel} bounds the open link moment, up to one normalizing constant, by \((1+b^{-6})\) times the constant of Siegel's step times \(\Sigma_{\varrho+1}\); this constant is the constant of Siegel's formula times the number of classes, the order of the reduction of the integral points of \(\Uc\), and \(\#\bar R_0\) (Definition~\ref{bp-def:bessel-rep}).  Apply Problem~\ref{pr:induction} with \(i=\varrho+1\).  The product of the \(c_i\), of the constant of Siegel's step and of \(c_{\mathrm{bot}}\) is at most \((1+b^{-9})^2(1+b^{-6})\,b^{-6k\#\Dc}\Dm\nu^2\), up to the same normalizing constant.  The factors are: \(1+b^{-9}\) for the constants, \((1+b^{-10})^f\le1+b^{-9}\) for the degrees (in the proof \(f\le b/1000\)), and \(1+b^{-6}\) from \(c_{\mathrm{bot}}\).  The error terms are at most \(b^{-10}\cdot b^{-6k\#\Dc}\Dm\nu^2\), because the large powers of \(b\) in \(c_{\mathrm{fr}}\) and \(c_{\mathrm{wt}}\) are paid by \(b^{11kf}\Di\le\Dm\) (Chapter~\ref{m1-chap}).  Since \(b\ge1000\), the following inequality collects these factors.

\begin{problem}\label{pr:constants}
For \(b\ge1000\): \((1+b^{-6})\bigl((1+b^{-9})^2(1+b^{-6})+b^{-10}\bigr)\le1+b^{-4}\).
\end{problem}

Problems~\ref{pr:last-entry}, \ref{pr:poisson-finite}, \ref{pr:schur}, \ref{pr:formal-degree-sum}, \ref{pr:induction} and~\ref{pr:constants} model the steps of the following theorem.  The proof proves it by these steps, with Theorems~\ref{thm:siegel}, \ref{thm:poisson} and~\ref{thm:sieve} and the statements read from the proof in this sheet.

\begin{theorem}\label{thm:open}
For the data, let \(\Dc\subseteq\cMc_0\) be a set of chosen places above \(p_0\), and for each \(w\in\Dc\) choose \(w\) or \(\bar w\) and a nonzero vector of the span of the frame modulo it, as in Definition~\ref{m2-def:fixer}; let the fixer of \(\Dc\) be the fixer of these vectors.  Then the link second moment of the terms \(\gamma\) with \(\disc(\gamma)\ne0\), over the reductions that are open and lie in the fixer of \(\Dc\), is at most
\[
 (1+b^{-4})\,b^{-6k\#\Dc}\,\Dm\,\nu^2 .
\]
\end{theorem}

\noindent\emph{Status}: read from the proof, where it is stated for an admissible filling over a CM field with Weil data, which the data provide.  The local bounds at nonzero frequencies and the matching of the constants are not derived here; see the remark below.

\paragraph{What has been proved.}  A term over an open reduction is a head, a last entry and a residual isometry (Section~\ref{sec:last}).  Its last entry has \(\lvert\lE\rvert^2-N^2=N\,h(y,y)\) (Problem~\ref{pr:last-entry}), so the residual isometries are counted by Siegel's formula (Theorem~\ref{thm:siegel}, modelled in Problems~\ref{pr:six-squares}--\ref{pr:densities}).  The heads are summed by Poisson summation one entry at a time: zero frequencies give integrals and the degrees \(d_\tau\), and the rank \(k\) pays for the rest (Problems~\ref{pr:fourier-five}--\ref{pr:formal-degree-sum}, Theorems~\ref{thm:poisson} and~\ref{thm:sieve}).  Taken as given: Theorems~\ref{thm:siegel}, \ref{thm:poisson} and~\ref{thm:sieve}, and the statements marked ``in the proof''.  Chapter~\ref{m4-chap} combines Theorem~\ref{thm:open} with the link bound and the cut of Chapter~\ref{m2-chap} (Theorems~\ref{m2-thm:link-bound} and~\ref{m2-thm:cut}), and bootstraps over fixers, to bound the second moment over all of \(\bar G\).

\begin{leanremark}
{\raggedright Theorem~\ref{thm:open} is \tl{New.Filling.linkSecondMoment\_isOpenElement\_le} in the Lean proof \cite{GL}.  Theorem~\ref{thm:siegel} is \tl{New.classSum\_eq\_mul\_normDensity}, Theorem~\ref{thm:poisson} is \tl{New.IsDualFamily.sum\_le\_of\_fourier}, and Theorem~\ref{thm:sieve} is read from \tl{New.Filling.sieve}.  Problem~\ref{pr:induction} is \tl{densitySum\_le}, Problem~\ref{pr:last-entry}, for \(V=Y_0\), is \tl{mul\_star\_dotProduct\_eq\_last\_frameProducts} with \glsrc{New/Spanning/MainTerm/GlobalToLocal.lean}{297}{the strict inequality}, Problem~\ref{pr:schur} is used as \tl{posSemidef\_prod\_mul}, and the matching of the constants is \tl{sieve\_constant\_le}.  The block as a sum over heads is \tl{exists\_amplitudeSum\_eq\_sum\_cosetAmplitude}, and the Kloosterman factorization is \tl{IsFillingData.cosetAmplitude\_eq\_kloostermanSum\_smul}.  The other statements marked ``in the proof'' are at these places of \cite{GL}, commit \texttt{72fead1}: \glsrc{New/Spanning/MainTerm/EnergyEstimate.lean}{90}{minors of open terms}, \glsrc{Lib/DiscreteSeries/Bruhat/BigCell.lean}{289}{the normal form} and \glsrc{Lib/DiscreteSeries/Bruhat/BigCell.lean}{90}{its minors}, \glsrc{New/Spanning/MainTerm/EnergyEstimate.lean}{666}{distinct blocks}, \glsrc{New/Spanning/MainTerm/EnergyEstimate.lean}{479}{Siegel's step for the blocks}, \glsrc{New/Spanning/MainTerm/SiegelTopBound.lean}{738}{Siegel's formula as an upper bound}, \glsrc{New/Spanning/MainTerm/SieveSystem.lean}{147}{the product of kernels}, \glsrc{New/Spanning/MainTerm/GlobalSieve.lean}{580}{nonnegative densities}, \glsrc{New/Spanning/MainTerm/SieveSystem.lean}{66}{the zero frequency}, \glsrc{New/Spanning/MainTerm/GlobalSieve.lean}{1177}{the decay}, \glsrc{New/Spanning/MainTerm/GlobalSieve.lean}{731}{the recurrence}, \glsrc{New/Spanning/Unfolding/Setup.lean}{113}{\(\varrho\ge10^5\)}, \glsrc{New/Spanning/Unfolding/Setup.lean}{128}{\(f\le b/1000\)}, and \glsrc{New/Spanning/Unfolding/Setup.lean}{515}{\(b^{11kf}\Di\le\Dm\)}.\par}
\end{leanremark}

\section*{Hints}

\subsection*{\S\ref{sec:last}}
\begin{description}
\item[\ref{pr:last-entry}] Expand \(h(\omega z,\omega z)=h(z,z)\), using \(y\perp z\).
\end{description}

\subsection*{\S\ref{sec:siegel}}
\begin{description}
\item[\ref{pr:six-squares}] In Sage or Python, enumerate the vectors of \(\Z^6\) with entries in \([-2,2]\).
\item[\ref{pr:local-counts}] By machine: convolve the histogram of the squares modulo \(m\) six times.  For the formula, detect the equation with \(\mathrm e(\cdot/p)\) and use the quadratic Gauss sum: every nonzero frequency gives the same sixth power.  Solutions modulo \(p\) with a unit coordinate lift by Hensel's lemma; compare the counts modulo \(3\) and \(9\).  Every solution has an odd coordinate; for \(e\ge3\), adding \(2^{e-1}\) to it changes the sum of squares by \(2^e\) modulo \(2^{e+1}\).
\item[\ref{pr:isotropic}] Induction on \(n\): \(x^{q+1}=c\) has \(q+1\) solutions in \(\mathbb F_{q^2}\) if \(c\in\mathbb F_q^\times\), one if \(c=0\).
\item[\ref{pr:real-density}] The unit ball of \(\R^6\) has volume \(\pi^3/6\).
\item[\ref{pr:densities}] \(\prod_{p\text{ odd}}(1-\chi_{-4}(p)p^{-3})=L(3,\chi_{-4})^{-1}\).  For \(j=5\), compute modulo \(25\) and \(125\) at \(5\), and at \(2\) use that \(y\mapsto(2+\ii)y\) is a bijection modulo \(2^e\) from the solutions of \(y^*y\equiv1\) to those of \(y^*y\equiv5\).
\end{description}

\subsection*{\S\ref{sec:poisson}}
\begin{description}
\item[\ref{pr:poisson-finite}] Fourier inversion and the triangle inequality.
\item[\ref{pr:kloosterman-sum}] Expand \(\lvert\mathrm{Kl}(j)\rvert^2\) and sum over \(j\) first.
\item[\ref{pr:three}] The rows of \(\mathfrak A_1\) are the characters of \(\Z/3\); the operator norm is multiplicative under tensor products.
\item[\ref{pr:schur}] The entrywise product is a principal submatrix of the Kronecker product; compare \cite[Thm.~7.5.3]{HJ}.
\item[\ref{pr:formal-degree-sum}] Apply the orthogonality relation with \(u_1=u_2=u_3=u_4=u\).
\item[\ref{pr:induction}] Induction on \(i\).
\item[\ref{pr:constants}] Each of \(b^{-6}\), \(b^{-9}\), \(b^{-10}\) is at most \(b^{-4}/100\).
\end{description}

\begin{thebibliography}{We65}
\bibitem[Br]{Br} H.~Braun, \emph{Zur Theorie der hermitischen Formen}, Abh.\ Math.\ Sem.\ Hansischen Univ.\ \textbf{14} (1941), 61--150.
\bibitem[GL]{GL} L.~Sernau, \emph{GRH for Dirichlet \(L\)-functions, in Lean~4}, a shorter rewriting of the Lean proof, text by Claude, \url{https://github.com/lsernau/grh-lean}, commit \texttt{72fead1}, readable edition \texttt{grh/}.
\bibitem[HJ]{HJ} R.~A.~Horn and C.~R.~Johnson, \emph{Matrix Analysis}, 2nd ed., Cambridge Univ.\ Press, 2013.
\bibitem[Kn]{Kn} A.~W.~Knapp, \emph{Representation Theory of Semisimple Groups: an Overview Based on Examples}, Princeton Univ.\ Press, 1986.
\bibitem[Si]{Si} C.~L.~Siegel, \emph{\"Uber die analytische Theorie der quadratischen Formen}, Ann.\ of Math.\ \textbf{36} (1935), 527--606.
\bibitem[Ta]{Ta} J.~T.~Tate, \emph{Fourier analysis in number fields and Hecke's zeta-functions}, in J.~W.~S.~Cassels and A.~Fr\"ohlich (eds.), \emph{Algebraic Number Theory}, Academic Press, 1967, Ch.~XV.
\bibitem[We65]{We65} A.~Weil, \emph{Sur la formule de Siegel dans la th\'eorie des groupes classiques}, Acta Math.\ \textbf{113} (1965), 1--87.
\end{thebibliography}

\end{document}
